\documentclass[troisieme]{fiche}
\metadonnees{10}{Le grand ménage}{Bloc B — Ce qu'on peut, ce qu'on doit}

\begin{document}
\entete

\objectifs{%
\textbf{Langue} : bilan des modaux ; expression écrite.\par
\textbf{Texte} : Kenneth Grahame, \emph{The Wind in the Willows} (1908), ch. \textsc{i}.\par
\textbf{Durée} : 1 h — exercices 20 min, texte et questions 20 min, version et expression
20 min.}

%% ===================================================================
\exercice[12 min]{Bilan des modaux}

\rappel{\textbf{Rappel.} Un modal dit ce que pense celui qui parle : c'est possible,
c'est obligatoire, c'est souhaitable. Il ne prend jamais de \emph{-s} et n'est jamais suivi
de \emph{to}.}

\consigne{Complétez avec \emph{can}, \emph{can't}, \emph{could}, \emph{must}, \emph{mustn't},
\emph{should}, \emph{had better} ou \emph{will}.}

\begin{anglais}
\begin{enumerate}[label=\arabic*.,itemsep=2.5mm,leftmargin=*]
  \item The Mole \trou{can't} stand the spring-cleaning any longer.
  \item ``I \trou{must} get out of this house,'' he said to himself.
  \item You \trou{should} rest a little: you have been working all morning.
  \item \trou{Can} you smell the spring in the air?
  \item He \trou{had better} put on his coat before he goes out.
  \item Moles \trou{can} dig very fast with their little paws.
  \item ``You \trou{mustn't} pass without paying,'' said the rabbit.
  \item One day he \trou{will} see the river for the first time.
  \item When he was younger he \trou{could} work all day without stopping.
  \item We \trou{should} not laugh at the rabbits.
\end{enumerate}
\end{anglais}

\note[Les trois emplois]{Capacité : 1, 4, 6, 9. Obligation ou interdiction : 2, 7. Conseil :
3, 5, 10. Avenir : 8. Un même modal change de valeur selon la phrase, et seul le contexte
tranche.}

%% ===================================================================
\exercice[8 min]{Expression écrite}
\consigne{Vous êtes la taupe. Le soir de cette journée, vous écrivez cinq ou six phrases dans
votre journal : ce que vous faisiez le matin, pourquoi vous êtes sorti, ce que vous avez
découvert. Employez au moins trois modaux différents. 80 mots.}

\lignes{10}

\corr{\textbf{Proposition.} \emph{March 20th.} I was spring-cleaning my little house this
morning. I had to sweep, and dust, and paint the walls with whitewash. At eleven o'clock I
couldn't stand it any longer. Something up above was calling me, so I threw down my brush and
ran up the tunnel. It was hard work, but I was able to get out at last. Now I am lying in the
warm grass of a great meadow. I must never spend another spring underground.}

\note[Barème]{Trois modaux différents, correctement construits (sans \emph{to}, sans
\emph{-s}) : 3 points. Récit au prétérit cohérent : 3 points. Quatre mots au moins pris dans
le lexique de la fiche : 2 points. Orthographe et ponctuation : 2 points.}

%% ===================================================================
\newpage
\section{Texte}

\noindent\small\textit{Une taupe fait son ménage de printemps dans sa maison souterraine.}\normalsize

\begin{texte}
The Mole had been working very hard all the morning, spring-cleaning his little home. First
with \gl[a broom]{brooms}{un balai}, then with \gl[a duster]{dusters}{un chiffon à
poussière}; then on ladders and steps and chairs, with a brush and a
\gl[a pail]{pail}{un seau} of \gl[whitewash]{whitewash}{du lait de chaux}; till he had dust
in his throat and eyes, and \gl[a splash]{splashes}{une éclaboussure} of whitewash all over
his black \gl[fur]{fur}{le pelage}, and an \gl[aching]{aching}{douloureux} back and weary
arms. Spring was moving in the air above and in the earth below and around him, penetrating
even his dark and \gl[lowly]{lowly}{humble} little house with its spirit of divine discontent
and \gl[a longing]{longing}{un désir ardent}. It was small wonder, then, that he suddenly
\gl[to fling, flung]{flung}{jeter} down his brush on the floor, said ``Bother!'' and ``O
blow!'' and also ``Hang spring-cleaning!'' and \gl[to bolt]{bolted}{détaler} out of the house
without even waiting to put on his coat. Something up above was calling him
\gl[imperiously]{imperiously}{impérieusement}, and he made for the steep little tunnel which
answered in his case to the \gl[a carriage-drive]{gravelled carriage-drive}{une allée
carrossable gravillonnée} owned by animals whose residences are nearer to the sun and air. So
he \gl[to scrape]{scraped}{gratter} and scratched and
\gl[to scrabble]{scrabbled}{gratter frénétiquement} and \gl[to scrooge]{scrooged}{se faufiler
en se tortillant} and then he scrooged again and scrabbled and scratched and scraped, working
busily with his little paws and muttering to himself, ``Up we go! Up we go!'' till at last,
pop! his \gl[a snout]{snout}{un museau} came out into the sunlight, and he found himself
rolling in the warm grass of a great meadow.

``This is fine!'' he said to himself. ``This is better than whitewashing!'' The sunshine
struck hot on his fur, soft \gl[a breeze]{breezes}{une brise} caressed his heated
\gl[a brow]{brow}{un front}, and after the \gl[seclusion]{seclusion}{le confinement} of the
\gl[cellarage]{cellarage}{les caves} he had lived in so long the \gl[a carol]{carol}{un
chant} of happy birds fell on his dulled hearing almost like a shout. Jumping off all his four
legs at once, in the joy of living and the delight of spring without its cleaning, he pursued
his way across the meadow till he reached the hedge on the further side.
\end{texte}
\source{Kenneth Grahame, \emph{The Wind in the Willows}, Londres, Methuen, 1908,
ch.~\textsc{i}.}

\partiel[12 min]{Questions}
\consigne{Répondez aux deux premières questions en français, aux deux dernières en anglais.}

\begin{enumerate}[label=\arabic*.,itemsep=2.5mm,leftmargin=*]
  \item Dans quel état la taupe se trouve-t-elle après sa matinée de ménage ?
    \lignes{2}
    \corr{De la poussière dans la gorge et dans les yeux, des éclaboussures de lait de chaux
    partout sur son pelage noir, le dos douloureux et les bras las.}
  \item Qu'est-ce qui la décide à tout abandonner ?
    \lignes{2}
    \corr{Le printemps, qui bouge dans l'air au-dessus et dans la terre autour d'elle, et qui
    pénètre jusque dans sa petite maison sombre. Quelque chose, là-haut, l'appelle
    impérieusement.}
  \item \en{``So he scraped and scratched and scrabbled and scrooged and then he scrooged
    again and scrabbled and scratched and scraped.'' Why does Grahame write the four verbs
    twice, the second time backwards?}
    \lignes{3}
    \corr{To make the reader feel how long and repetitive the climb is: the same movements,
    over and over. Reading the sentence takes almost as long as the tunnel does, and the
    words come back in reverse as if he were retracing his own passage.}
  \item \en{Find in the text three words that describe noises or movements of digging.}
    \lignes{2}
    \corr{\emph{scraped}, \emph{scratched}, \emph{scrabbled}, \emph{scrooged}, et
    \emph{pop!} au moment où il sort. Tous commencent par le même groupe de consonnes, qui
    imite le bruit.}
\end{enumerate}

%% ===================================================================
\section{Version}
\consigne{Traduisez le passage en français. Il suit immédiatement le texte : la taupe arrive à
la haie.}

\begin{version}
``Hold up!'' said an elderly rabbit at the \gl[a gap]{gap}{une trouée, un passage}. ``Sixpence
for the privilege of passing by the private road!'' He was \gl[to bowl over]{bowled over}{
renverser} in an instant by the impatient and \gl[contemptuous]{contemptuous}{méprisant}
Mole, who \gl[to trot]{trotted}{trotter} along the side of the hedge
\gl[to chaff]{chaffing}{railler} the other rabbits as they \gl[to peep]{peeped}{jeter un coup
d'œil} hurriedly from their holes to see what the \gl[a row]{row}{un vacarme} was about.
``Onion-sauce! Onion-sauce!'' he remarked \gl[jeeringly]{jeeringly}{d'un ton narquois}, and
was gone before they could think of a thoroughly satisfactory reply.
\end{version}
\source{\emph{Ibid.}}

\lignes{12}

\corr{\textbf{Proposition de traduction.} \og Halte-là ! \fg{} dit un lapin âgé, posté dans la
trouée. \og Six pence pour avoir le privilège d'emprunter le chemin privé ! \fg{} Il fut
renversé en un instant par la taupe, impatiente et méprisante, qui trotta le long de la haie
en raillant les autres lapins, tandis qu'ils sortaient précipitamment le nez de leurs terriers
pour voir d'où venait ce vacarme. \og Sauce à l'oignon ! Sauce à l'oignon ! \fg{} lança-t-elle
d'un ton narquois, et elle avait disparu avant qu'ils pussent trouver une réponse
pleinement satisfaisante.}

\note[Choix de traduction 1]{\emph{He was bowled over by the Mole} : passif avec complément
d'agent. Ici le français garde le passif sans peine — c'est quand l'agent manque qu'il faut
passer par \og on \fg.}

\note[Choix de traduction 2]{\emph{Onion-sauce!} : en Angleterre, le lapin se mange avec une
sauce à l'oignon. La taupe leur rappelle donc, en passant, ce qui les attend. Une note est
indispensable : traduire seul ne transmet rien.}

\note[Choix de traduction 3]{\emph{was gone} : \emph{be} + participe passé exprime ici un
état, non une action — \og elle avait disparu \fg, et non \og elle fut partie \fg.}

%% ===================================================================
\partiel{Lexique à mémoriser}
\begin{multicols}{2}\small
\begin{itemize}[label=--,itemsep=0.8mm,leftmargin=*]
  \item \textbf{a mole} — une taupe
  \item \textbf{spring-cleaning} — le ménage de printemps
  \item \textbf{a broom} — un balai
  \item \textbf{dust} — la poussière \textit{(fiche 4)}
  \item \textbf{a throat} — une gorge
  \item \textbf{weary} — las
  \item \textbf{fur} — le pelage, la fourrure
  \item \textbf{a ladder} — une échelle
  \item \textbf{steep} — raide \textit{(fiche 3)}
  \item \textbf{a meadow} — un pré \textit{(fiche 3)}
  \item \textbf{a hole} — un trou
  \item \textbf{a reply} — une réponse
\end{itemize}
\end{multicols}

\end{document}
